\section*{الفصل \ref{ch:b3:extension-strike-slip} --- التكتونية الاستطالية والانزياحية}

\begin{solution}{exo:b3:extension-strike-slip:1}
تنحدر الطبقات بمقدار $50 - 35 = \ang{15}$، بعيدًا عن الصدوع؛ $\beta =
\sin\ang{50}/\sin\ang{35} = 0.766/0.574 = 1.34$.
\end{solution}

\begin{solution}{exo:b3:extension-strike-slip:2}
التدرّج الأيمن نحو اليمين: حوض. التدرّج الأيمن نحو اليسار: ظهرة. التدرّج الأيسر
نحو اليسار: حوض. التدرّج الأيسر نحو اليمين: ظهرة.
\end{solution}

\begin{solution}{exo:b3:extension-strike-slip:3}
$105/20 = \qty{5.25}{km/Myr} = \qty{5.25}{mm/yr}$، قريب من
\qty{6}{mm/yr}: لا يلزم أن يكون المعدّل ثابتًا خلال \qty{20}{Myr}.
\end{solution}

\begin{solution}{exo:b3:extension-strike-slip:4}
\qty{200}{\percent} تعني $\Delta L = 2L_0$: $L_0 = \qty{117.5}{km}$ و$L
= \qty{352.5}{km}$. في الشمال، $\Delta L = 0.5L_0$: $L_0 = \qty{470}{km}$
و$L = \qty{705}{km}$.
\end{solution}

\begin{solution}{exo:b3:extension-strike-slip:5}
$563/20 = 28$ إلى $563/15 = \qty{38}{mm/yr}$. تمثّل الإزاحة حدًّا أدنى، والانجراف المقيس بالمسح عبر نطاق الصدع يتضمّن أيضًا
الحركةَ على الصدوع المجاورة في المنظومة
والانفعالَ المرن المخزَّن
بين الزلازل.
\end{solution}

\begin{solution}{exo:b3:extension-strike-slip:6}
$t = 90 - \text{الانحدار}$. عند \ang{60}: $(1 + 0.6\times1.732)/(1 - 0.6\times
0.577) = 2.039/0.654 = 3.12$؛ عند \ang{50}: $1.715/0.497 = 3.45$؛ عند
\ang{40}: $1.504/0.285 = 5.28$. يحدث الانغلاق عند $\arctan0.6 = \ang{31.0}$.
\end{solution}

\begin{solution}{exo:b3:extension-strike-slip:7}
$\arctan0.4 = \ang{21.8}$ و$\arctan0.2 = \ang{11.3}$. الصدع المنحدر
بمقدار \ang{15} لا ينزلق إلا إذا كان $\mu < \tan\ang{15} = 0.27$.
\end{solution}

\begin{solution}{exo:b3:extension-strike-slip:8}
$2.5\times4.0 = \qty{10.0}{km^2}$ و$2.5\times(4.0 - 1.5) =
\qty{6.25}{km^2}$.
\end{solution}

\begin{solution}{exo:b3:extension-strike-slip:9}
$132/1.19 = \qty{110.9}{km}$، تمدّدت بمقدار \qty{21.1}{km}، وهذا يستغرق
$21.1/6 = \qty{3.5}{Myr}$.
\end{solution}

\begin{solution}{exo:b3:extension-strike-slip:10}
$3.0/\sin\ang{60} = \qty{3.464}{km}$ قبل و$3.0/\sin\ang{45} =
\qty{4.243}{km}$ بعد: \qty{0.779}{km} لكل كتلة ($\beta = 1.22$). تنحدر الطبقات
بمقدار $60 - 45 = \ang{15}$. تدور كل كتلة كوحدة واحدة: الدوران الذي
يخفّض انحدار الصدوع المنحدرة، لنفرض نحو الشرق، يدير خطًّا أفقيًّا في
الكتلة بحيث يرتفع نحو الشرق، وتنحدر الطبقات نحو الغرب، عكس اتجاه
الصدوع.
\end{solution}

\begin{solution}{exo:b3:extension-strike-slip:11}
في نطاق أيسر شرقي--غربي يتّجه $\sigma_1$ بزاوية \ang{045}. قصوص R عند
$\phi/2 = \ang{20}$ عن النطاق على جانب $\sigma_1$: الامتداد \ang{070}،
أيسر. قصوص R$'$ عند $90 - 20 = \ang{70}$: الامتداد \ang{020}، أيمن. الكسور الشدّية
الموازية للإجهاد $\sigma_1$: الامتداد \ang{045}.
\end{solution}

\begin{solution}{exo:b3:extension-strike-slip:12}
عند $\mu = 0.6$ تتشكّل الصدوع عند \ang{60.5} وتنغلق عند \ang{31.0}: $\beta =
\sin\ang{60.5}/\sin\ang{31.0} = 0.870/0.515 = 1.69$. عند $\mu = 0.85$: من
\ang{65.2} إلى \ang{40.4}، $\beta = 0.908/0.648 = 1.40$. الصدوع ذات الاحتكاك الأدنى
تدور أكثر قبل الانغلاق: تمتدّ الصخور الضعيفة أكثر قبل أن
تتشكّل صدوع جديدة.
\end{solution}

\begin{solution}{pb:b3:extension-strike-slip:1}
\textbf{1.} كانت الطبقات أفقية ودارت مع الكتل؛
والصدوع دارت بالزاوية نفسها، فيكون $\theta_0 = 42 + 18 = \ang{60}$،
زاوية انحدار أندرسون عند $\phi = \ang{30}$.

\textbf{2.} $\beta = \sin\ang{60}/\sin\ang{42} = 0.8660/0.6691 = 1.294$.

\textbf{3.} $24.0/1.294 = \qty{18.5}{km}$؛ امتدّت المنطقة بمقدار
\qty{5.5}{km} (\qty{5.46}{km}).

\textbf{4.} $t = \ang{48}$: $(1 + 0.6\times0.900)/(1 - 0.6\times1.111) =
1.540/0.334 = 4.62$، مقابل $3.12$ للصدوع الجديدة عند \ang{60.5}، أي
بزيادة النصف.

\textbf{5.} عند $\arctan0.6 = \ang{31.0}$؛ $\beta = \sin\ang{60}/\sin\ang{31}
= 1.68$ على الأكثر. قبل ذلك، ستقطع صدوع جديدة عند نحو \ang{60} الكتل،
وسيتولّى جيلٌ ثانٍ من \omterm{def:b3:extension-strike-slip:domino}{التصدّع الدومينوي} الأمر.

\textbf{6.} يهبط كل عمود من الجدار المعلّق بمقدار هبوط الصدع الواقع
تحته مقسومًا على الإزاحة الأفقية؛ يكون الصدع أشدّ انحدارًا قرب السطح،
ويتفلطح بالابتعاد عنه، لذلك تهبط الأعمدة القريبة من الصدع أكثر ما تهبط والبعيدة
أقلّ: تميل الطبقات نحو الصدع.

\textbf{7.} $D = A/H = 16.0/3.2 = \qty{5.0}{km}$.

\textbf{8.} $D$ يتناسب طرديًّا مع $A$: $\qty{5.0}{km} \pm \qty{0.5}{km}$.

\textbf{9.} $3.2\times(5.0 - 2.0) = \qty{9.6}{km^2}$.

\textbf{10.} الإزاحة الأفقية، \qty{3.2}{km}.

\textbf{11.} يزداد سمك الحشوة نحو الصدع على شكل إسفين، وتكون طبقاته الأقدم
أشدّ ميلًا من الأحدث: ترسّبت كل طبقة بينما
\omterm{def:b3:extension-strike-slip:detachment}{الطية المحدبة الانقلابية} والحيّز المجاور للصدع تزدادان اتّساعًا.

\textbf{12.} $z = P/(\rho g) = \num{5.5e8}/(2800\times9.8) = \qty{20.0}{km}$.

\textbf{13.} $20.0/(24 - 16) = \qty{2.5}{km/Myr} = \qty{2.5}{mm/yr}$.

\textbf{14.} $20.0/\sin\ang{30} = \qty{40}{km}$ من الانزلاق، و$40\cos\ang{30} =
\qty{34.6}{km}$ من الاستطالة الأفقية.

\textbf{15.} لا: \ang{15} أدنى من انحدار القفل البالغ \ang{31}. إمّا أن الصدع
انزلق حين كان أشدّ انحدارًا ثم دُوِّر (\omterm{def:b3:extension-strike-slip:core-complex}{المفصل المتدحرج})، أو
كان ضعيفًا، مع $\mu < 0.27$ بفعل الطين أو السوائل عند ضغط مرتفع.

\textbf{16.} بمقدار $45 - 15 = \ang{30}$: مع ابتعاد الصفيحة العليا، يرتفع
الجدار المرتكز غير المحمَّل بالتوازن الأيزوستاتي وينثني، فيدور الصدع عبر
المفصل.

\textbf{17.} \qty{34.6}{km} مقابل \qty{8.7}{km} لنصف الغرابن (السؤال
23): أي أربعة أضعاف. تقع المعقّدات اللبّية المتحوّلة حيث تكون الاستطالة أكبر، كما في
وسط إقليم الحوض والسلاسل.

\textbf{18.} صدع أيسر ينتقل يسارًا: \omterm{def:b3:extension-strike-slip:bends}{انحناءة محرِّرة}، فينفتح حوض.

\textbf{19.} $\qty{5}{mm/yr}\times\qty{1.2}{Myr} = \qty{6.0}{km}$.

\textbf{20.} $4.0\times6.0 = \qty{24}{km^2}$؛ $3.0 + 6.0 = \qty{9.0}{km}$ طولًا؛
$9.0/4.0 = 2.25$.

\textbf{21.} $132/10 = 13$ إلى $132/7 = 19$: حوض البحر الميّت أكثر
استطالة بكثير، إذ تتجاوز إزاحته عرضه بمرّات عديدة.

\textbf{22.} \omterm{def:b3:extension-strike-slip:bends}{انحناءة مقيِّدة}: ظهرة مدفوعة إلى الأعلى بين الفروع؛ في
المقطع، \omterm{def:b3:extension-strike-slip:flower}{بنية زهرية} موجبة، ينفرج الصدع نحو الأعلى إلى فروع
ذات مركّبة عكسية، مثل التي يتراوح عرضها بين 100 و\qty{300}{m} على
صدع البحر الميّت الانزياحي.

\textbf{23.} $3.2 + 5.46 = \qty{8.66}{km}$: نصف غرابن ميرو قد استُطيل
بمقدار \textbf{\qty[text-series-to-math=true]{8.7}{km}}.
\end{solution}
